
\section{Estándares de ingeniería}

\subsection{Estándares de Redes y Comunicaciones:}

\begin{itemize}
  \item \textbf{IEEE 802.11}: El estándar define las especificaciones técnicas y operativas para las redes inalámbricas locales (WLAN). Incluye protocolos que aseguran la interoperabilidad entre dispositivos, la transmisión segura de datos, la autenticación, y la gestión del espectro radioeléctrico. Este estándar es fundamental para cualquier dispositivo que se conecte a redes Wi-Fi, asegurando la comunicación efectiva y segura en dichas redes.

  \item \textbf{IEEE 802.15}: El conjunto de estándares IEEE 802.
\end{itemize}



\subsection{Estándares de Software y Seguridad de la Información:}

\begin{itemize}
  \item \textbf{ISO/IEC 12207:} El estándar proporciona un marco para los procesos del ciclo de vida del software, abarcando desde la concepción hasta la retirada del software. Define procesos clave como la adquisición, la especificación de requisitos, el diseño, la implementación, la verificación, la validación y el mantenimiento del software. Este estándar es esencial para garantizar una gestión de calidad y una metodología estructurada en el desarrollo y mantenimiento del software en dispositivos vestibles y otros sistemas electrónicos. 

  \item \textbf{ISO/IEC 27001:} Especifica los requisitos para establecer, implementar, operar, monitorizar, revisar, mantener y mejorar un Sistema de Gestión de Seguridad de la Información (SGSI). Este estándar es crucial para garantizar la confidencialidad, integridad y disponibilidad de la información, especialmente en dispositivos vestibles que recopilan y transmiten datos sensibles. 

\end{itemize}

\subsection{Estándares de Usabilidad y Accesibilidad:}

\begin{itemize}
  \item \textbf{WCAG 2.1:} Las Pautas de Accesibilidad para Contenido Web (WCAG) 2.1 proporcionan un conjunto de recomendaciones para hacer el contenido web más accesible, especialmente para personas con discapacidades. Abarca aspectos como la percepción, la operabilidad, la comprensibilidad, y la robustez del contenido web. Aunque se centra en la web, las pautas también pueden aplicarse a aplicaciones móviles y otros medios digitales, asegurando que las interfaces de usuario sean inclusivas y accesibles para todos los usuarios, independientemente de sus capacidades. 

\end{itemize}

\subsection{Estándares de Interoperabilidad y Conformidad:}

\begin{itemize}
  \item \textbf{ASTM F2761:} Pr

  \item \textbf{ISO/IEC 30118:} spositivos conectados. 
\end{itemize}

\subsection{Estándares de Arquitectura y Evaluación de Calidad:}

\begin{itemize}
  \item \textbf{IEEE P2413:} El estándar proporciona un marco arquitectónico para el diseño y la implementación de sistemas y soluciones basadas en Internet de las Cosas (IoT). Establece definiciones, principios arquitectónicos y modelos para facilitar la interoperabilidad, la seguridad y la privacidad en entornos IoT. Es crucial para asegurar una arquitectura coherente y bien estructurada en proyectos que involucren dispositivos vestibles conectados a redes IoT. 

\end{itemize}
